\chapter{Lepton Masses and Hierarchy}
\label{ch:leptons}

\begin{quote}
\textit{This chapter establishes the EDC ontology for charged leptons:
electron, muon, and tau as brane-localized excitations in a mode tower,
with the mass hierarchy arising from mode index.}
\end{quote}

% ════════════════════════════════════════════════════════════════════════════════
\section{The Lepton Mass Puzzle}
\label{sec:lepton_puzzle}
% ════════════════════════════════════════════════════════════════════════════════

\begin{tcolorbox}[edcGuardrail, title=\textbf{Epistemic Status}]
\textbf{What is established:}
\begin{itemize}[nosep]
    \item Mode tower structure with $n = 0, 1, 2$ for $(e, \mu, \tau)$ \tagI{}
    \item Brane-dominant ontology for charged leptons \tagP{}
    \item Electron stability as ground-state consequence \tagDc{}
\end{itemize}

\textbf{What is NOT derived:}
\begin{itemize}[nosep]
    \item Explicit mass values from 5D action
    \item Mass ratio formula
    \item Why these particular masses
\end{itemize}
\end{tcolorbox}

The charged lepton masses span a factor of $\sim 3500$ \tagBL{}:
\begin{center}
\begin{tabular}{lcc}
\toprule
\textbf{Lepton} & \textbf{Mass (MeV)} & \textbf{Ratio to $m_e$} \\
\midrule
Electron ($e$) & $0.511$ & $1$ \\
Muon ($\mu$) & $105.7$ & $207$ \\
Tau ($\tau$) & $1777$ & $3477$ \\
\bottomrule
\end{tabular}
\end{center}

The Standard Model takes these as input parameters. EDC proposes a geometric origin.

% ════════════════════════════════════════════════════════════════════════════════
\section{Physical Picture: Brane-Dominant Excitations}
\label{sec:lepton_physical}
% ════════════════════════════════════════════════════════════════════════════════

\begin{mechanism}{Charged Leptons as Mode Tower}
\textbf{Step 1: Charged leptons are brane-localized.}
Unlike the proton (a bulk junction) or neutron (excited junction), charged
leptons are \textbf{brane-dominant excitations}---they live entirely within
the brane layer, localized on the observer-facing side \tagP{}.

\textbf{Step 2: The brane supports a mode tower.}
The thick-brane potential $V(\xi)$ supports discrete bound states. These are
labeled by mode index $n = 0, 1, 2, \ldots$, with energy increasing with $n$ \tagDc{}.

\textbf{Step 3: Charged leptons are the first three modes.}
\begin{itemize}[nosep]
    \item $n = 0$: Electron (ground state)---stable
    \item $n = 1$: Muon (first excited state)---decays to electron
    \item $n = 2$: Tau (second excited state)---decays to muon or electron
\end{itemize}

\textbf{Step 4: Mass hierarchy from mode energy.}
Higher modes have more nodes in their wavefunction and higher energy:
$E_n > E_{n-1}$. The mass hierarchy $m_\tau > m_\mu > m_e$ follows directly \tagI{}.

\textbf{Step 5: Stability is spectral.}
The electron is stable because it is the ground state---there is no lower-energy
charged mode to decay into \tagDc{}.
\end{mechanism}

% ════════════════════════════════════════════════════════════════════════════════
\section{Electron: The Ground State}
\label{sec:electron}
% ════════════════════════════════════════════════════════════════════════════════

\begin{tcolorbox}[colback=green!5!white, colframe=green!50!black,
    title=\textbf{Electron Ontology}]
The electron is a \textbf{stable topological defect} localized on the
observer-facing layer of the brane.

\textbf{Key properties:}
\begin{enumerate}[nosep]
    \item \textbf{Localization:} Confined to $y \approx +\delta/2$ (observer-facing side)
    \item \textbf{Mode index:} $n = 0$ (ground state of charged sector)
    \item \textbf{Stability:} No lower-lying charged mode exists---decay forbidden
    \item \textbf{Mass:} $m_e = 0.511$ MeV (baseline, not derived)
\end{enumerate}
\end{tcolorbox}

\textbf{Why the electron is the universal decay endpoint:}

All weak decay chains involving charged leptons terminate at the electron
because it is the ground state of the charged brane-mode spectrum. This is
not an assumption---it is a spectral consequence \tagDc{}.

\begin{center}
\begin{tabular}{ll}
\toprule
\textbf{Observable} & \textbf{EDC Explanation} \\
\midrule
$\tau_e > 10^{28}$ years & Ground state: no lower mode to decay to \\
$m_e$ smallest charged mass & Mode $n=0$ has lowest energy \\
$e^-$ in all leptonic decays & Only allowed output for charged channel \\
\bottomrule
\end{tabular}
\end{center}

% ════════════════════════════════════════════════════════════════════════════════
\section{Muon: First Excited State}
\label{sec:muon}
% ════════════════════════════════════════════════════════════════════════════════

\begin{tcolorbox}[colback=blue!5!white, colframe=blue!50!black,
    title=\textbf{Muon at a Glance}]
\textbf{Baseline} \tagBL{}:
\begin{itemize}[nosep]
    \item Decay: $\mu^- \to e^- + \bar\nu_e + \nu_\mu$ with $\tau_\mu = 2.197 \times 10^{-6}$ s
    \item Purely leptonic (no hadrons)
    \item Energy: $m_\mu c^2 = 105.7$ MeV available
\end{itemize}

\textbf{EDC View} \tagP{}:
\begin{itemize}[nosep]
    \item Muon = brane-dominant excitation (mode $n = 1$)
    \item Decay is internal mode relaxation within brane layer
    \item No bulk$\to$brane pumping---energy already in brane
\end{itemize}
\end{tcolorbox}

\textbf{Muon decay as mode relaxation:}

\begin{enumerate}[nosep]
    \item \textbf{5D cause:} The muon occupies an unstable mode ($n=1$) within the brane layer.
    \item \textbf{Brane response:} Mode redistribution occurs entirely within the brane---no bulk junction relaxation.
    \item \textbf{3D output:} The frozen projection emits $e^-$, $\nu_\mu$, $\bar{\nu}_e$.
\end{enumerate}

This is a ``clean room'' test of the EDC pipeline: same mechanism as neutron decay
but without bulk-core complications \tagDc{}.

% ════════════════════════════════════════════════════════════════════════════════
\section{Tau: Second Excited State}
\label{sec:tau}
% ════════════════════════════════════════════════════════════════════════════════

\begin{tcolorbox}[colback=purple!5!white, colframe=purple!50!black,
    title=\textbf{Tau at a Glance}]
\textbf{Baseline} \tagBL{}:
\begin{itemize}[nosep]
    \item Mass: $m_\tau = 1777$ MeV
    \item Lifetime: $\tau_\tau = 2.9 \times 10^{-13}$ s (much shorter than muon)
    \item Decay channels: Leptonic ($\sim 35\%$) and hadronic ($\sim 65\%$)
\end{itemize}

\textbf{EDC View} \tagP{}:
\begin{itemize}[nosep]
    \item Tau = brane-dominant excitation (mode $n = 2$)
    \item Higher mode $\Rightarrow$ more decay channels available
    \item Hadronic decays: $m_\tau > m_\pi + m_\nu$ opens pion channels
\end{itemize}
\end{tcolorbox}

\textbf{Why tau has more decay channels:}

The tau mass exceeds the pion mass, allowing hadronic decay modes that are
kinematically forbidden for the muon. This is purely kinematic \tagBL{}, but
in EDC it reflects the higher-mode structure having more overlap with
hadronic brane excitations \tagI{}.

% ════════════════════════════════════════════════════════════════════════════════
\section{Mass Hierarchy from Mode Index}
\label{sec:mass_hierarchy}
% ════════════════════════════════════════════════════════════════════════════════

The mass hierarchy $m_e \ll m_\mu \ll m_\tau$ maps to mode indices:

\begin{center}
\begin{tabular}{lccc}
\toprule
\textbf{Lepton} & \textbf{Mode $n$} & \textbf{Mass (MeV)} & \textbf{Status} \\
\midrule
Electron & $0$ & $0.511$ & Ground state (stable) \\
Muon & $1$ & $105.7$ & First excited (decays) \\
Tau & $2$ & $1777$ & Second excited (decays) \\
\bottomrule
\end{tabular}
\end{center}

\textbf{Mass ratio pattern:}
\begin{equation}
    \frac{m_\mu}{m_e} \approx 207, \qquad \frac{m_\tau}{m_\mu} \approx 17
    \quad \tagBL{}
\end{equation}

The Koide formula provides an empirical fit \tagBL{}:
\begin{equation}
    \frac{m_e + m_\mu + m_\tau}{(\sqrt{m_e} + \sqrt{m_\mu} + \sqrt{m_\tau})^2} = \frac{2}{3}
    \quad \text{(within } 0.02\%\text{)}
\end{equation}

\textbf{Status:} The mode tower structure \tagI{} explains why there are three
charged leptons with hierarchical masses. The specific mass values require
deriving the brane potential $V(\xi)$ from first principles \tagOpen{}.

% BACKFILL-GAP11: Yukawa from overlap BEGIN
\subsection{Mass from Overlap Integrals (Not Yukawa)}

In the Standard Model, fermion masses arise from Yukawa couplings:
$m_f = y_f v / \sqrt{2}$ where $y_f$ is a free parameter \tagBL{}.

In EDC, fermion masses arise from 5D wavefunction overlaps at the brane:

\begin{mechanism}{Fermion Mass from Localization Overlap}
\textbf{5D setup} \tagDc{}:
A fermion mode $f_n(\xi)$ has energy $E_n$ from the brane potential.
The effective 4D mass depends on the \textbf{overlap integral}:
\begin{equation}
m_f^{(n)} \propto \int_{-\delta/2}^{+\delta/2} |f_n(\xi)|^4 \, d\xi \equiv I_4^{(n)}
\label{eq:overlap_mass}
\end{equation}

\textbf{Physical meaning:}
\begin{itemize}[nosep]
\item Higher modes ($n > 0$) have more nodes $\Rightarrow$ smaller $I_4$
\item Ground state ($n = 0$) has largest overlap $\Rightarrow$ different coupling
\item Mass hierarchy emerges from mode structure, not free Yukawa couplings
\end{itemize}

\textbf{What this replaces:} The SM Yukawa parameters $y_e, y_\mu, y_\tau$
become derived quantities from mode eigenvalues and overlaps \tagDc{}.
\end{mechanism}

\begin{tcolorbox}[colback=yellow!5!white, colframe=yellow!50!black,
    title=\textbf{Dictionary: EDC Overlap $\leftrightarrow$ SM Yukawa}]
\textbf{SM description} \tagBL{}: Yukawa coupling $y_f$ is a free parameter.

\textbf{EDC description} \tagDc{}: $y_f^{\text{(eff)}} \propto I_4^{(n)}$
is determined by mode geometry.

\textbf{Status:} This identification is structural \tagI{}. Numerical verification
requires solving the BVP for explicit $f_n(\xi)$ profiles (\OPR{21}).
\end{tcolorbox}
% BACKFILL-GAP11: Yukawa from overlap END

% ════════════════════════════════════════════════════════════════════════════════
\section{Connection to Generation Structure}
\label{sec:lepton_generations}
% ════════════════════════════════════════════════════════════════════════════════

The mode index $n = 0, 1, 2$ corresponds to the $\mathbb{Z}_3$ generation
structure from Chapter~\ref{ch:generations}:

\begin{itemize}[nosep]
    \item Each charged lepton generation occupies one of three angular sectors
    \item The neutrino partner occupies the same sector (see Chapter~\ref{ch:neutrinos})
    \item Mass hierarchy arises from mode energy, not angular position
\end{itemize}

\textbf{Key distinction:}
\begin{itemize}[nosep]
    \item \textbf{Generation number} (which family): from $\mathbb{Z}_3$ angular sector
    \item \textbf{Mass value} (how heavy): from mode eigenvalue in brane potential
\end{itemize}

% ════════════════════════════════════════════════════════════════════════════════
\section{Falsifiability}
\label{sec:lepton_falsify}
% ════════════════════════════════════════════════════════════════════════════════

\begin{tcolorbox}[edcCanonical, title=\textbf{What Would Falsify the Mode Tower Picture}]
\begin{enumerate}
\item \textbf{Electron decay observed:}
      Would mean ground state is not stable---mode structure fails.

\item \textbf{Lighter charged particle discovered:}
      Electron would not be mode $n=0$.

\item \textbf{4th charged lepton:}
      Would require mode $n=3$, inconsistent with $\mathbb{Z}_3$ structure.

\item \textbf{Mass hierarchy inverted} (e.g., $m_e > m_\mu$):
      Would contradict mode energy ordering.
\end{enumerate}
\end{tcolorbox}

% ════════════════════════════════════════════════════════════════════════════════
\section{Chapter Summary}
\label{sec:lepton_summary}
% ════════════════════════════════════════════════════════════════════════════════

\begin{statusbox}{Key Takeaways}
\begin{enumerate}
\item \textbf{Charged leptons are brane-dominant excitations:}
      They live entirely within the brane layer, unlike bulk junctions (proton) \tagP{}.

\item \textbf{Mode tower structure:}
      $n = 0$ (electron), $n = 1$ (muon), $n = 2$ (tau) \tagI{}.

\item \textbf{Electron stability:}
      Ground state has no lower mode to decay to \tagDc{}.

\item \textbf{Muon/tau decay:}
      Internal mode relaxation within brane, same pipeline as neutron \tagDc{}.

\item \textbf{Mass hierarchy:}
      Higher modes have higher energy $\Rightarrow$ $m_\tau > m_\mu > m_e$ \tagI{}.
\end{enumerate}
\end{statusbox}

\begin{tcolorbox}[colback=orange!5, colframe=orange!50!black,
    title=\textbf{Epistemic Audit: Chapter 7}]
\begin{center}
\begin{tabular}{lll}
\toprule
\textbf{Claim} & \textbf{Status} & \textbf{Source} \\
\midrule
$m_e, m_\mu, m_\tau$ values & \tagBL{} & PDG \\
Brane-dominant ontology & \tagP{} & Postulated \\
Mode tower $n = 0, 1, 2$ & \tagI{} & Identified with masses \\
Electron stability & \tagDc{} & Ground-state consequence \\
Muon decay mechanism & \tagDc{}/\tagP{} & Pipeline from Ch. 1 \\
Koide formula match & \tagBL{}/\tagI{} & Empirical pattern \\
\bottomrule
\end{tabular}
\end{center}
\end{tcolorbox}

\textbf{Next:} Chapter~\ref{ch:generations} explains why exactly three generations
exist, and Chapter~\ref{ch:neutrinos} treats the neutrino partners as edge modes.

